\subsection{停用之后还能怀上吗：各类避孕方法的生育力恢复}

"\textbf{用了这么多年避孕药，以后会不会怀不上？}"——这是门诊与咨询里出现频率最高的疑问之一。直接给答案：\textbf{现有的避孕方法都不会造成不孕}；停用后的生育力取决于年龄与原有的身体状况，而不是"用了多久""用了多少"。但不同方法的恢复速度差别很大，这一点常被完全忽略——有人在停药两个月后就开始备孕，有人选用了恢复期最长的注射剂，却以为"停药就能怀"。

\textbf{一、复方口服避孕药、避孕贴片与阴道环}

\begin{itemize}
\item \textbf{恢复时间}：停药后排卵多在 1--3 个月内恢复，不少人第一个周期即可恢复排卵——\textbf{不需要"停药三个月到半年再怀孕"}，这个流传很广的说法没有依据；
\item \textbf{长期使用的影响}：多项大样本随访的结论是，长期服用复方口服避孕药不降低停药后的受孕概率。也就是说，服药 5 年与服药 1 年，停药后的生育力没有可测量的差别；
\item \textbf{一个需要提前知道的常见情形}：停药后可能出现短暂闭经，尤其是原本周期就不规律的人。若超过 3 个月没有月经，应就诊排查其他原因（多囊卵巢综合征、甲状腺问题、高催乳素血症等），而不是笼统归因于"避孕药把身体吃坏了"；
\item \textbf{服药期间怀孕的风险}：如按规则服用，概率很低。若停药后随即备孕，叶酸可以在计划怀孕前 3 个月开始补充。
\end{itemize}

\textbf{二、单纯孕激素避孕药}

恢复速度与复方制剂相近，停药后即可备孕，无需要等待期。

\textbf{三、DMPA（每三个月一次的注射剂）——恢复最慢的一种}

\begin{itemize}
\item 药物在体内持续释放，\textbf{停药后平均需 6--9 个月恢复排卵}，少数人可延长至 12--18 个月；
\item 因此若计划在 1--2 年内怀孕，DMPA 通常不是首选；已经使用者若决定生育，应尽早停用并与医生沟通;
\item 一个值得知道的补充：停用后怀孕的时间虽晚，但\textbf{有过 DMPA 使用史并不影响最终能否怀孕}；延迟的是"何时开始"，不是"能不能"。
\end{itemize}

\textbf{四、皮下埋植与宫内节育器（含激素 IUS）}

\begin{itemize}
\item \textbf{取出后恢复最快}：激素 IUS 的孕激素作用以局部为主、清除快；皮下埋植取出后激素水平迅速回落。二者取出后\textbf{即可恢复排卵}，当月或下一个周期怀孕都是常见的;
\item \textbf{含铜宫内节育器}：取出即恢复，无延迟。放置期间或取出后的生育问题，多与既往盆腔感染对输卵管的影响相关，而与节育器本身无关;
\item 若取出后已尝试超过 1 年（35 岁以下）或 6 个月（35 岁以上）仍未怀孕，应按不孕症的常规路径检查。
\end{itemize}

\textbf{五、绝育手术：性质完全不同}

必须把两类问题分开：

\begin{itemize}
\item 前面几种方法的讨论前提都是"停用后会恢复"；
\item \textbf{绝育手术是永久性的}。它不恢复，只能通过复通手术争取，而复通后的妊娠率明显低于未绝育人群（详见本书绝育手术与"绝育后悔"相关小节）。因此绝育前的知情评估极为重要——不要把它当成"反正以后能接通"的选项。
\end{itemize}

\textbf{六、什么时候应该就医排查}

\begin{itemize}
\item 35 岁以下，未避孕尝试 12 个月未孕；35 岁及以上，尝试 6 个月未孕；
\item 停用任何避孕方法后超过 3--6 个月（DMPA 使用者按延迟规律相应放宽）仍无月经且无排卵迹象；
\item 有月经长期不规律、子宫内膜异位症、既往盆腔炎、腹部或盆腔手术史者，宜更早评估，不必等到上述时间点。
\end{itemize}

\textbf{七、为什么这一节值得写}

把"避孕"和"生育力"的关系讲清楚，可以减少两类真实的伤害：

\begin{itemize}
\item 一是\textbf{因为怕不孕而不敢用有效避孕方法}——改用低效方法或不用方法，最终以意外怀孕收场；
\item 二是\textbf{停用后短期未孕就陷入恐慌}，去接受不必要的检查与治疗，或把它解释为对过去某种选择的"惩罚"。
\end{itemize}

\textit{【编者注】}影响生育力的主要变量，是年龄、既往盆腔感染与手术史、子宫内膜异位症、抽烟、体重与内分泌状况，而不是"曾经用过哪种避孕方法"。计划怀孕前做一次孕前评估，比纠结过去的避孕史更有价值。

